% !PRINTED 551
\weberpage{551}
% running head: \S.\,151. Ideale. 551
%<<CONT
mit der Zahl überein, die bei \gesperrt{Dedekind} die \gesperrt{Norm des Ideals}
heisst. Sie soll also auch hier so genannt werden, und wir setzen
demnach, wenn das Functional $\varphi$ zu dem Ideal $\mathfrak{a}$ gehört,

\weq{2}{N_a(\varphi) = N(\mathfrak{a}).}

Die Norm eines Ideals ist also immer eine natürliche Zahl.

Ist $\mathfrak{p}$ ein Primideal und $p$ die durch $\mathfrak{p}$ theilbare natürliche
Primzahl, so ist

\weq{3}{N(\mathfrak{p}) = p^f,}

und $f$ heisst der \gesperrt{Grad} des Primideals $\mathfrak{p}$.

In den Congruenzen (\S.\,149, 150) können nach diesen Festsetzungen
 statt der Functionale auch die entsprechenden Ideale
als Moduln angesehen werden.

Es sei schliesslich noch ein Wort über die \gesperrt{Kummer'sche}
Schöpfung der idealen Zahlen beigefügt\footnote{Vgl. \gesperrt{Dirichlet-Dedekind}, Zahlentheorie, \S.\,176.}.

Sind $\alpha, \beta$ irgend zwei Zahlen aus $\mathfrak{o}$, die den grössten
gemeinschaftlichen Theiler $\chi$ haben, und ist $\beta = \chi\varphi$, so wird
die Congruenz

\weq{4}{\alpha\omega \equiv 0 \pmod{\beta}}

nur für solche Zahlen $\omega$ aus $\mathfrak{o}$ befriedigt sein, die durch $\varphi$
theilbar sind. Das Functional $\varphi$ wird aber im Allgemeinen
keiner Zahl associirt sein. Statt nun die Functionale zu benutzen,
kann man sich ein neues Begriffssystem schaffen, in dem
man die sämmtlichen Zahlen von $\mathfrak{o}$ zu Paaren zusammenfasst
$(\alpha, \beta)$, und kann diese Paare \glqq \gesperrt{ideale Zahlen}\grqq{} in $\mathfrak{o}$ nennen;
genau in derselben Weise, wie man durch Paarung der ganzen
rationalen Zahlen die Brüche, oder der reellen Zahlen die
complexen Zahlen gebildet hat.

Man nennt dann die Zahl $\omega$ durch die ideale Zahl $(\alpha, \beta)$
theilbar, wenn sie der Congruenz (1) genügt. % [sic] printed reference reads \textquotedbl{}(1)\textquotedbl{}; the congruence introduced on this page is numbered (4)
Dies ist der Standpunkt
 der \gesperrt{Kummer'schen} Theorie. Damit aber diese Definition
der Theilbarkeit fruchtbar sei, muss noch festgestellt werden,
unter welchen Voraussetzungen man zwei ideale Zahlen mit verschiedenen
 Elementen, $(\alpha, \beta), (\alpha', \beta')$, als gleich zu betrachten
hat, wie sich die wirklichen Zahlen in $\mathfrak{o}$ unter die idealen Zahlen
einordnen, endlich was man unter dem Product zweier idealen
Zahlen zu verstehen hat, damit eine durch eine ideale Zahl
%CONT>>
